\documentclass[12pt]{article}

% --- Packages ---
\usepackage[a4paper,margin=1in]{geometry}
\usepackage{microtype}
\usepackage{enumitem}
\usepackage{siunitx}
\usepackage{array}
\usepackage{float}
\usepackage[hidelinks]{hyperref}

% --- Settings ---
\sisetup{detect-all}
\setlist[itemize]{topsep=4pt,itemsep=2pt,parsep=0pt}
\setcounter{secnumdepth}{3}

% --- Metadata ---
\title{BMS Firmware Architecture Overview}
\author{Kaushik Vada}
\date{\today}

\begin{document}
\maketitle

\section{Firmware Architecture Overview}

The Battery Management System (BMS) firmware is designed to ensure safe operation, continuous monitoring, and protection of the lithium-ion battery pack across multiple runtime conditions. The firmware executes within an RTOS environment to support deterministic task scheduling and timely fault-response behavior.

\subsection{BMS Modes of Operation}

The BMS operates in two primary power modes:

\begin{itemize}
    \item \textbf{SHIP Mode:} Low-power hibernation state intended for storage and shipping conditions. Only essential safety monitoring remains active to minimize power consumption and prevent unnecessary battery drain.
    \item \textbf{NORMAL Mode:} Full operational state in which sensing, protection, balancing, and communication functions are active. The BMS continuously monitors pack health and reports data to the host system.
\end{itemize}

\subsection{RTOS State Machines}

Under NORMAL Mode, core safety tasks execute periodically according to their timing requirements.

\subsubsection{Temperature Monitoring Task (ADC Sampler)}

Responsible for thermal fault protection:

\begin{itemize}
    \item Periodically samples all thermistors using the ADC
    \item Compares each reading against configured over-temperature thresholds
    \item On fault detection:
    \begin{itemize}
        \item Activate audible buzzer
        \item Flash the LED corresponding to the overheated cell (one LED per thermistor)
        \item Transmit alert and cell identifier to the host GUI for visualization
    \end{itemize}
\end{itemize}

This task runs at a high frequency due to the sudden nature of thermal runaway events.

\subsubsection{Overcurrent Monitor Task}

Provides high-speed protection against current surges:

\begin{itemize}
    \item Samples pack current via the Coulomb Counter register in the BMS IC
    \item Compares the value against software-defined overcurrent thresholds
    \item On violation:
    \begin{itemize}
        \item Trigger buzzer and fault indication LEDs
        \item Send alert to host GUI for user awareness
    \end{itemize}
\end{itemize}

Like temperature monitoring, this task requires rapid execution to catch transient fault events.

\subsubsection{BMS Handler Task}

Overall system management and slow-timescale functions:

\begin{itemize}
    \item Schedules and executes cell balancing operations
    \item Maintains overall state machine control and internal flags
    \item Periodically transmits telemetry and status to the GUI
    \item Operates based on RTOS timer interrupts or asynchronous safety triggers
\end{itemize}

Because voltage imbalance progresses slowly, this task operates at a lower rate.

\subsection{System Summary}

Table~\ref{tab:task_rates} summarizes major tasks and their roles in ensuring battery safety and long-term operation.

\begin{table}[h!]
\centering
\begin{tabular}{|c|c|p{7cm}|}
\hline
\textbf{Task} & \textbf{Execution Rate} & \textbf{Primary Purpose} \\
\hline
Temperature Monitoring & 10--100\,\si{ms} & Detect overheating early and prevent thermal runaway \\
\hline
Overcurrent Monitor & 10--100\,\si{ms} & Detect short circuits / excessive discharge currents \\
\hline
BMS Handler & 1--10\,\si{s} & Manage cell balancing and maintain long-term pack health \\
\hline
\end{tabular}
\caption{RTOS task execution rates and associated safety roles.}
\label{tab:task_rates}
\end{table}

In summary, the BMS firmware leverages its RTOS-based architecture to prioritize rapid fault detection while sustaining overall system health. Audible, visual, and digital alert paths provide both local and remote visibility into safety-critical events.

\end{document}